\begin{textsample}{POS Dim 4 – 2022\_c6 – Score 3.00 – 2022\_c6/enade\_2022\_dir\_7d44b6aec0.txt}  \label{ex:f4_pos_089}
CAPÍTULO 7 ANÁLISE TÉCNICA DA PROVA Este capítulo tem por objetivo apresentar o desempenho dos estudantes concluintes de Direito no Enade/2022.
Para isso, foram calculadas as estatísticas básicas da prova em seu todo ( Subseção 7.1.1 ), bem como as estatísticas relacionadas aos Componentes de Formação Geral ( Subseção 7.1.2 ) e de Conhecimento Específico da área ( Subseção 7.1.3 ).
Dadas as suas características, foram analisadas, em separado, as questões objetivas ( Seção 7.2 ) e as questões discursivas ( Seção 7.3 ).
Tomando -se como base as duas questões discursivas do Componente de Formação geral, nas Subseções 7.3.1.6 a 7.3.1.8, são apresentados os resultados e os comentários sobre a correção das respostas em relação à Língua Portuguesa.
Para as questões objetivas, foram disponibilizados os índices de Facilidade e de Discriminação Ponto-Bisserial, também em separado, para os Componentes de Formação Geral ( Subseção 7.2.1 ) e de Conhecimento Específico ( Subseção 7.2.2 ).
De cada componente, uma das questões foi escolhida para exemplificar a análise gráfica, relacionando -se as alternativas escolhidas pelos estudantes ( inclusive o gabarito ) com o número de acertos no componente.
No Anexo IV, apresenta -se a íntegra da análise gráfica para todas as questões objetivas.
Para cada uma das questões discursivas, os conteúdos dos tipos mais comuns de respostas dos estudantes são apresentados e comparados com o padrão de resposta esperado ( ver Anexo IX ).
Nas tabelas, constam as seguintes estatísticas das notas¹⁰ : média do desempenho na prova, erro-padrão da média, desvio-padrão, nota mínima, mediana e nota máxima para cada um de seus componentes.
Tais estatísticas contemplam o total de estudantes concluintes da Área de Direito inscritos e presentes à prova do Enade/2022, tendo em vista agregações, ou por grandes regiões e o país como um todo, ou por Categoria administrativa, organização acadêmica da IES e modalidade de oferta dos cursos.
Em relação aos gráficos de distribuição de notas, o intervalo considerado foi de 10 unidades, aberto à esquerda e fechado à direita, com exceção do primeiro intervalo, [ 0 ; 10 ], fechado em ambos os extremos.
Para os gráficos de distribuição das notas das questões discursivas, foram consideradas mais duas categorias : questão em branco¹¹ e nota zero. 10 Uma definição dessas estatísticas pode ser encontrada no Glossário inserido no final desse relatório. 11 Nesse grupo estão incluídas também as respostas classificadas como nulas ou desconsideradas. 7.1 ESTATÍSTICAS BÁSICAS DA PROVA Esta seção apresenta estatísticas selecionadas e histogramas da nota geral ( 7.1.1 ) e de cada componente : Formação Geral ( 7.1.2 ) e Conhecimento Específico ( 7.1.3 ).
São também apresentadas estatísticas selecionadas de subpopulações, caracterizadas por grande região, categoria administrativa, organização acadêmica e modalidade de oferta. 7.1.1 ESTATÍSTICAS BÁSICAS GERAIS Na Tabela 7.1, são apresentadas as estatísticas básicas da prova, por grande região, dos estudantes concluintes de Direito.
Como se pode verificar, a partir dos dados mostrados na Tabela 7.1, a média das notas da prova, como um todo, foram 45,2 ( nas seções seguintes serão analisados os Componentes de Formação Geral e de Conhecimento Específico ).
Os estudantes da região Norte obtiveram a média mais baixa ( 43,6 ), e os da região Sul alcançaram a média mais alta ( 46,9 ).
As demais médias foram : 45,7, na região Nordeste, 44,9, na região Sudeste e 44,2, na região Centro-Oeste.
O desvio-padrão para o Brasil foi de 15, sendo o maior desvio-padrão encontrado na região Nordeste ( 15,2 ), e o menor, na região Sul ( 14,6 ), indicando uma dispersão um pouco menor das notas desta última região.
A região que obteve a maior nota máxima foi a região Sul ( 92,4 ), ao passo que a região que atingiu a menor nota máxima foi a Norte ( 86,5 ).
A Mediana do Brasil, como um todo, foram 44,9, sendo a maior Mediana obtida na região Sul ( 47,1 ), e a menor obtida na região Norte ( 43 ).
Todas as regiões obtiveram nota mínima zero.
O comportamento das notas dos estudantes de todo o Brasil pode ser observado no Gráfico 7.1, que apresenta um histograma com a distribuição das mesmas : uma distribuição unimodal com a moda no intervalo ( 30 ; 40 ].
Na Tabela 7.2, são apresentadas informações referentes à Média da nota final, desagregadas por Categoria administrativa, organização acadêmica e modalidade de oferta dos cursos.
Os dados da Tabela 7.2 mostram que, em relação à Categoria administrativa, os estudantes das IES públicas obtiveram Média mais alta ( 53,9 ) que a Média das IES privadas ( 43,8 ).
O Desvio-padrão para as IES públicas foi de 15,2, indicando uma dispersão das notas nesta Categoria administrativa.
No tocante à organização acadêmica, os estudantes de CEFET/IF obtiveram a média de 50,5, enquanto os de universidades, centros universitários e faculdades obtiveram, respectivamente, média de 46,9, 44,5 e 43. 7.1.2 ESTATÍSTICAS BÁSICAS NO COMPONENTE DE FORMAÇÃO GERAL Na Tabela 7.3, são apresentadas as estatísticas básicas em relação ao componente da prova que avalia a Formação Geral dos estudantes concluintes.
A nota de Formação Geral é obtida levando -se em conta os acertos nas oito questões objetivas e a nota média obtida nas questões discursivas 1 e 2, sendo a nota de cada questão discursiva obtida levando -se em conta 80 \% da nota relativa ao conteúdo e 20 \% da nota relativa ao desempenho linguístico.
Conforme os dados da Tabela 7.3, os estudantes de todo o Brasil obtiveram, no Componente de Formação Geral da prova, desempenho médio de 57,6.
Quanto à variabilidade, o desvio-padrão das notas dos estudantes do Brasil foi 18,4.
A maior média foi obtida na região Sul ( 59,2 ) e a menor, na região Centro-Oeste ( 55,7 ).
As demais médias foram : 58,7 na região Nordeste, 57,3 na região Sudeste e 56 na região Norte.
Já o maior desvio-padrão foi obtido na região Centro-Oeste ( 18,9 ), e o menor, na região Sul ( 18 ).
Os demais desvios foram : 18,5 na região Norte, 18,5 na região Nordeste e 18,3 na região Sudeste.
A maior nota no Componente de Formação Geral da prova do Enade/2022 foi 100, obtida por, pelo menos, um estudante nas quatro regiões : Sul, Sudeste, Nordeste e Centro-Oeste.
A menor nota máxima foi obtida na região Norte ( 98,6 ).
A Mediana do Brasil, como um todo, foram 59,1 ; sendo a menor Mediana encontrada na região Centro-Oeste ( 56,9 ), e a maior, na região Sul ( 60,8 ).
A nota mínima nessa parte foi zero em todas as regiões.
No Gráfico 7.2, é possível verificar a avaliação do desempenho dos estudantes do Componente de Formação geral, a partir do histograma da distribuição das notas correspondentes.
Nesse gráfico, a distribuição é unimodal, com moda em [ 60 ; 70 ].
Na Tabela 7.4, são apresentadas as informações referentes ao desempenho dos concluintes do Componente de Formação geral, em diferentes agregações : categoria administrativa, organização acadêmica e modalidade de oferta.
A partir dos dados da Tabela 7.4, observa -se que as médias dos tipos de organização acadêmica foram : 59,9 para as universidades ; 56,7 para os centros universitários ; 54,7 para as faculdades e 68,8 para os CEFET/IF. 7.1.3 ESTATÍSTICAS BÁSICAS DO COMPONENTE DE CONHECIMENTO ESPECÍFICO Na Tabela 7.5, são apresentadas as estatísticas básicas referentes ao Componente de Conhecimento Específico da área de Direito.
A nota do componente de Conhecimentos Específicos leva em conta as notas da parte objetiva e a média das notas obtidas nas três questões discursivas da parte discursiva da área.
De acordo com a Tabela 7.5, a média do desempenho dos estudantes do Brasil, para o componente específico da prova, foi 41.
A maior média foi obtida na região Sul ( 42,8 ), e a menor na região Norte 39,4.
As demais médias foram : 41,3 na região Nordeste e 40,7 na região Sudeste.
Quanto à variabilidade das notas, o desvio-padrão do Brasil foi 15,9 ; sendo o maior desvio-padrão observado na região Nordeste ( 16,2 ) e o menor, nas regiões Norte e Sul ( 15,5 e 15,6 ).
Os demais desvios foram : 15,7 na região Centro-Oeste e 15,9 na região Sudeste.
A Mediana das notas dos estudantes de todo o Brasil foi 40,6.
A maior mediana ocorreu na região Sul ( 42,7 ), e a menor, na região Norte ( 38,7 ).
As demais medianas foram : 39,7 na região Centro-Oeste, 40,8 na região Nordeste e 40,1 na região Sudeste.
A nota máxima do Brasil, como um todo, foram 91,6 sendo obtida por pelo menos um estudante na região Sul.
As demais notas máximas foram : 89,6 na região Nordeste, 88,1 na região Centro-Oeste ; 90,6 na região Sudeste ; 86,3 na região Norte.
A maior nota mínima foi zero em todas as regiões.
Assim como nos Gráficos 7.1 e 7.2, antes mostrados, no Gráfico 7.3, apresentado adiante, é possível ser feita uma avaliação do desempenho de concluintes de Direito, em relação ao Componente de Conhecimento Específico, com um histograma da distribuição das notas correspondentes.
Esta também é uma distribuição unimodal, e o grupo modal é o ( 30 ; 40 ], o mesmo do grupo modal para a Formação Geral e para a prova, de forma geral.
Na Tabela 7.6, é apresentada uma comparação dos resultados em relação a categoria administrativa, a organização acadêmica e a modalidade de oferta, levando -se em conta o desempenho de estudantes do Componente de Conhecimento Específico da prova.
De acordo com os dados expostos na Tabela 7.6, no que se refere à organização acadêmica, a maior média foi a pública ( 49,3 ), vindo a seguir a dos CEFET/IF ( 44,4 ), em seguida as universidades ( 42,6 ), seguida pela dos centros universitários ( 40,5 ) e, por fim, pelas faculdades ( 39,1 ).
O maior desvio-padrão foi das universidades e públicas ( 16,6 ).
As universidades públicas obtiveram a maior nota máxima ( 91,6 ).
Os centros universitários obtiveram nota Máxima 89,6 ; as faculdades, 89,2 e os CEFET/IF, 71,3.
As Medianas foram : 42,5 nas Universidades, 40,1 nos centros universitários, 38,6 nas faculdades e 42,9 nos CEFET/IF.
A maior nota mínima foi 23, nos CEFET/IF.
Quanto à categoria administrativa, observa -se um comportamento semelhante àquele da parte de Formação Geral e à prova, como um todo. 7.2 ANÁLISE DAS QUESTÕES OBJETIVAS Esta seção apresenta estatísticas selecionadas e histogramas das questões objetivas de Formação Geral ( 7.2.1 ) e de Conhecimento Específico ( 7.2.2 ).
São, também, apresentadas e comparadas as médias das subpopulações caracterizadas por grande região.12 12 As fórmulas utilizadas para o coeficiente de correlação ponto-bisserial e para o índice de facilidade estão apresentadas no Capítulo 2. 7.2.1 COMPONENTE DE FORMAÇÃO GERAL Na Tabela 7.7, são apresentadas as estatísticas básicas relativas às oito questões objetivas do Componente da prova que abrange a Formação Geral dos estudantes.
De acordo com os dados da Tabela 7.7, em relação ao Componente de Formação geral, a Média do Brasil foi 58,7.
A menor Média foi encontrada na região Norte ( 56,4 ), e a maior, na região Sul ( 60,9 ).
As demais médias foram : 59,8 na região Nordeste, 58,3 na região Sudeste e 57 na região Centro-Oeste.
O desvio-padrão do Brasil foi 22,4 ; sendo o maior desvio-padrão encontrado na região Centro-Oeste ( 22,7 ), e o menor desvio-padrão encontrado na região Sul ( 21,9 ).
As medianas do Brasil, como um todo, e das cinco regiões ( Sul, Nordeste, Norte, Sudeste e Centro-Oeste ) foram iguais a 62,5.
A nota máxima ( 100,0 ) foi alcançada nas cinco grandes regiões.
As notas mínimas ( 0,0 ) foram iguais para todas as regiões.
Na Tabela 7.8, exposta mais adiante, são apresentados o Índice de Facilidade e o Índice de Discriminação ( Ponto-Bisserial ) para cada uma das questões objetivas do Componente de Formação geral.
De acordo com os dados da Tabela 7.8, as questões objetivas do Componente de Formação Geral, segundo o Índice de Facilidade, foram assim avaliadas : nenhuma teve o Índice de Facilidade classificado como Muito Fácil.
Quatro questões foram classificadas como fácil, com valor de 0,71 a 0,82.
Três questões foram classificadas como médio por terem valores de 0,41, 0,42 e 0,55.
Oito questões foram classificadas como muito bom, com valores entre 0,42 e 0,53.
Enquanto uma questão foi classificada na categoria Difícil, com valor de 0,34.
O Índice de Facilidade variou de 0,34 a 0,82 e o de Discriminação, de 0,42 a 0,53.
Na Tabela 7.9, é apresentada a distribuição das questões, levando -se em conta, simultaneamente, a classificação dos dois índices.
Os dados da Tabela 7.9 atestam que oito questões que tiveram Índice de Discriminação Muito bom figuraram em três níveis de dificuldade – Fácil, Médio e Difícil : quatro classificadas na categoria Fácil, três classificadas na categoria Médio e uma na categoria Difícil.
No Gráfico 7.4, cada uma das cinco curvas representa o percentual de respostas em determinada alternativa da questão, em função do número de acertos dos estudantes nessa parte da prova ( Formação geral/Múltipla Escolha ), antes de possíveis eliminações pelo critério Ponto-Bisserial. 7.2.2 COMPONENTE ESPECÍFICO No Gráfico 7.4, cada uma das cinco curvas representa o percentual de respostas em determinada alternativa da questão, em função do número de acertos dos estudantes nessa parte da prova ( Formação geral/Múltipla Escolha ), antes de possíveis eliminações pelo critério Ponto-Bisserial.
Os gráficos relativos às demais questões de Formação Geral constam do Anexo IV.
Na Tabela 7.10, são apresentadas as estatísticas básicas em relação às questões objetivas do Componente de Conhecimento Específico da prova de Direito por grande região.
Como se verifica a partir dos dados da Tabela 7.10, a Média do Brasil para o Componente de Conhecimento Específico da Prova foi 44,8.
A menor Média foi observada na região Norte ( 43,1 ), e a maior, na região Sul ( 46,7 ).
O desvio-padrão de todo o Brasil foi 17,6, sendo o maior desvio-padrão encontrado na região Nordeste ( 18 ), e o menor, na região Norte ( 17,2 ).
A mediana de todo o Brasil foi 45,8, o mesmo valor da mediana encontrada nas regiões, Nordeste, Sul e Sudeste.
Na região Centro-Oeste, a mediana foi 41,7.
A nota máxima da prova foi 100, obtida nas questões objetivas do Componente de Conhecimento Específico, por, pelo menos, um estudante das regiões Sudeste e Sul.
Nas regiões Centro-Oeste e Nordeste a nota máxima foi 95,8 e 91,7 na região Norte.
A nota mínima foi zero, obtida por pelo menos um estudante em todas as cinco regiões.
Na Tabela 7.11, são apresentados os Índices de Facilidade e de Discriminação ( Ponto-Bisserial ) das questões objetivas do Componente de Conhecimento Específico, para os estudantes de Direito.
Para facilitar a diferenciação das questões, serão usadas as mesmas cores da Tabela 7.8, anteriormente apresentada, para as diferentes classificações dos Índices de Facilidade e de Discriminação.
A partir do Índice de Facilidade obtido, apresentado na Tabela 7.11, pode -se concluir que a maioria das questões objetivas da prova foi considerada de nível médio ( 12 ).
Onze questões foram consideradas Bom, 3 questões foram consideradas fácil, 1 questão foi considerada muito difícil, 3 como fraco, 8 como muito bom e 10 como difícil.
O Índice de Facilidade variou de 0,11 a 0,80 e o de Discriminação, de – 0,01 a 0,48.
Na Tabela 7.12, é apresentada a distribuição das questões levando -se em conta, simultaneamente, a classificação dos dois índices.
Os dados da Tabela 7.12 atestam que onze questões que tiveram Índice de Discriminação “ Bom ” figuraram em três níveis de dificuldade – Fácil, Médio e Difícil : uma classificada na categoria Fácil, cinco classificadas na categoria Médio e cinco na categoria Difícil.
Oito questões que tiveram Índice de Discriminação “ Muito Bom ” figuraram em dois níveis de dificuldade – Fácil e Médio : seis classificadas na categoria Médio e duas na categoria Fácil.
Cinco questões que tiveram Índice de Discriminação médio figuraram em dois níveis de dificuldade – Médio e Difícil : uma classificada na categoria Médio e quatro na categoria Difícil.
Três questões que tiveram Índice de Discriminação “ Fraco ” figuraram em dois níveis de dificuldade – Difícil e Muito Difícil : duas classificada na categoria Difícil e uma, na categoria Muito Difícil.
A título de exemplo das análises do comportamento das questões objetivas, no Gráfico 7.5, analisa -se a Questão 33 do Componente de Conhecimento Específico.
Como revela o Gráfico 7.11, a questão 24 do Componente de Conhecimento Específico foi considerada, pelas respostas dos estudantes avaliados na prova, como uma questão com nível de Facilidade Médio, 0,54, ou seja, 54 \% dos estudantes assinalaram, acertadamente, a opção C correspondente ao gabarito.
Como já comentado, seu Índice de Discriminação foi igual a 0,45, classificado como Muito Bom.
No mesmo gráfico, cada uma das cinco curvas representa o percentual de respostas em determinada alternativa da Questão 24, em função do número de acertos dos estudantes nessa parte da prova, antes de possíveis eliminações de questões pelo critério Ponto-Bisserial.
A alternativa correta C, representada no gráfico pela curva em verde, foi escolhida em maiores proporções pelos estudantes com desempenho melhor no Componente de Conhecimento Específico.
Já as alternativas incorretas, também denominadas distratores, foram selecionadas, principalmente, por aqueles com notas mais baixas. 7.3 ANÁLISE DAS QUESTÕES DISCURSIVAS Esta seção apresenta estatísticas selecionadas e histogramas das questões discursivas de Formação Geral ( 7.3.1 ) e de Componente Específico ( 7.3.2 ).
Expõe e compara, também, as médias de subpopulações, caracterizadas por grande região.
Cumpre notar que nem todas as questões passam pelo crivo dos corretores.
As respostas depois de digitalizadas passam por um pré-processamento de identificação automático.
Algumas destas questões são definidas como um “ branco automático ” ( que não são enviadas para os professores para correção, a não ser que na segunda fase, a de inspeção visual, seja constatada uma resposta ).
O critério usado neste pré-processamento, é verificar para cada questão, a partir da quantidade de pixels, o que provavelmente, é uma resposta em branco.
A partir da separação destes documentos que foram considerados em branco, eles são enviados para uma verificação visual, por uma equipe de colaboradores, devidamente capacitados.
Deste modo, quaisquer erros, deste pré-processamento, são encontrados, e as questões preenchidas são enviadas para a correção dos professores.
Eventualmente, algumas respostas em branco não são detectadas nesse pré-processamento automático.
A presença de pequenos riscos, ou manchas, pode fazer o programa não identificar uma questão como um possível branco.
Estas são enviadas para a correção, mas são classificadas como “ branco ” pelos professores, quando da correção. 7.3.1 COMPONENTE DE FORMAÇÃO GERAL As análises dos resultados de desempenho dos estudantes de Direito, nas duas questões discursivas relativas à Formação geral, encontram -se na Tabela 7.13 e no Gráfico 7.6, cujos dados são analisados de forma sequenciada.
Na seção 7.3.1.7 deste relatório na seção 7.3.1.7 deste relatório.
A nota de cada questão é composta de 80 \% da nota da correção de conteúdo e 20 \% da nota de desempenho linguístico.
Na Tabela 7.13, observa -se que os estudantes de todo o Brasil obteve, em Formação geral, média 56 nas questões discursivas.
Pode -se notar, também, que o desvio-padrão nesse conjunto de questões foi de 23.
A maior média foi obtida na região Nordeste ( 56,9 ) e a menor, na região Centro-Oeste ( 53,7 ).
A Mediana de todo o Brasil, neste componente, foi 60,5 e a maior Mediana foi a da região Nordeste, de 61,8.
A nota máxima ( 100 ) foi obtida nas cinco regiões.
A nota mínima ( 0,0 ) foi a mesma em todas as regiões do Brasil.
No Gráfico 7.6, está representada a distribuição das notas nas questões discursivas do Componente de Formação geral.
Como mostra o Gráfico 7.6, a moda da distribuição ocorre no intervalo [ 60 ; 70 ], com frequência de aproximadamente 22,5 \% dos participantes, sendo que, no intervalo [ 0 ; 10 ], inclui -se, além da nota zero, a frequência de estudantes que deixaram esse tipo de questão em branco.
A distribuição é sempre crescente, exceto do sétimo intervalo, ( 60 ; 70 ], para frente.
Na sequência, os resultados verificados para cada uma das questões discursivas de Formação Geral serão apresentados, estabelecendo -se relações com os conteúdos abordados em cada uma delas.
Os comentários das Bancas de docentes corretores, a respeito do observado na correção das respostas dos estudantes, suas impressões e conclusões serão apresentados junto à análise de cada questão.
Cumpre esclarecer que, tendo em vista que as questões discursivas de Formação Geral são padronizadas, ou seja, constam de todas as provas, os comentários da Banca são os mesmos para todas as carreiras acadêmicas, sendo direcionados a todos os estudantes que participaram do Enade/2022.
A seguir, serão analisados o desempenho linguístico dos estudantes da Área de Direito nas duas questões discursivas de Formação Geral do Enade/2022, apresentando -se os resultados obtidos e os comentários da banca de correção para cada questão. 7.3.1.1 ANÁLISE DE CONTEÚDO DA QUESTÃO DISCURSIVA 1 DO COMPONENTE DE FORMAÇÃO GERAL Os dados de Direito, obtidos a partir das respostas à Questão 1 do Componente de Formação geral, encontram -se na Tabela 7.14 e no Gráfico 7.7, expostos após a apresentação dos dados da referida tabela.
Verifica -se, a partir dos dados da Tabela 7.14, que na Questão 1 do Componente de Formação Geral – de desempenho abaixo ao da outra questão do mesmo componente –, os estudantes de todo o Brasil tiveram média 53,9.
A maior média para a Questão 1 foi obtida na região Nordeste ( 55 ), e a menor, na região Centro-Oeste ( 51,7 ).
Quanto à variabilidade das notas, o desvio-padrão de todo o Brasil foi 25,7.
O menor desvio-padrão foi obtido na região Sudeste ( 25,3 ), e o maior desvio-padrão foi obtido pela região Centro-Oeste ( 26,5 ).
A Mediana do Brasil, como um todo, foram 55,5, sendo igual no Norte, Sudeste e no Centro-Oeste.
A nota máxima da Questão Discursiva 1 foi 100,0, obtida por, pelo menos, um estudante nas cinco regiões.
Já a nota mínima para todas as regiões foi zero.
No Gráfico 7.7, apresenta -se a distribuição das notas na Questão Discursiva 1 do Componente de Formação geral.
O Gráfico 7.7 mostra que mais de 5 \% dos participantes deixaram em branco a resposta à questão discursiva 1.
A moda dessa distribuição ocorre no intervalo ( 50 ; 60 ] com uma frequência maior que 20 \% dos participantes. 7.3.1.2 COMENTÁRIOS SOBRE A CORREÇÃO DE CONTEÚDO DAS RESPOSTAS À QUESTÃO DISCURSIVA 1 A Questão 1 abordou um tema atual, com dois itens para os estudantes responderem : a ) “ discuta por que os investimentos em educação, ciência, tecnologia e inovação são estratégicos e contribuem para o desenvolvimento de um país ” ; b ) “ explique como o fomento público ao desenvolvimento científico pode atender à justiça social e à equidade, em contextos como o da pandemia de Covid-19 ”.
O tema contempla como os investimentos em educação, ciência, tecnologia e inovação ( ECTI ) cooperam para um país se desenvolver e como o fomento público, vinculado ao desenvolvimento científico, pode atender à justiça social, à equidade e à população em momentos de pandemia.
Esse tema é atual, pois o mundo vivenciou recentemente uma pandemia que, somente no Brasil, deixou mais de 700 mil mortos, sendo exaustivamente comentada na televisão e nos meios digitais.
Logo, almeja -se que os respondentes justifiquem a necessidade de investimentos em ECTI com foco em questões como : investimentos público em P\&D por meio de universidades públicas ; pesquisas realizadas em outros países que possibilitaram o acesso mais célere às vacinas e serviram de modelo para o que foi desenvolvido no Brasil ; pesquisadores de universidades, tanto \textbf{brasileiras} quanto estrangeiras como protagonistas para orientar e informar a respeito da Covid-19 por meio dos canais de comunicação ; desenvolvimento de equipamentos de segurança, como máscara de nariz ou de rosto, respiradores, que foram implementados e viabilizados por pesquisa e extensão em universidades e escolas técnicas.
Esperava -se que os candidatos soubessem o básico a respeito do tema, pois o assunto é bastante atual.
Quanto ao item B, esperava -se que os candidatos abordassem questões como : o papel das universidades e dos centros de pesquisa para a justiça e a equidade ; o investimento público em pesquisa, tecnologia e educação para garantir acesso igualitário da população aos recursos de saúde, podendo fazer um contraponto com os investimentos privados ; a política de distribuição de vacinas gratuitas por meio do Sistema Único de Saúde, que possibilitou o acesso às vacinas no Brasil ; e as questões econômicas mundiais que determinaram uma distribuição desigual de recursos de saúde em escala global, com os países mais ricos retendo a maior parte desses recursos.
Atualmente, é comprovado que os países dominantes são aqueles que se desenvolvem e dominam as ferramentas tecnológicas.
As pesquisas mostram que os países que, mesmo em crise, investiram em ciência, tecnologia e inovação são os que conseguiram se reerguer mais rapidamente e sobressair -se aos demais países em crises pandêmicas.
Portanto, a questão foi elaborada de maneira oportuna e adequada ao universo do estudante universitário brasileiro.
No item A, o estudante poderia focar em diversos argumentos, entre eles : como ocorreu o investimento público em pesquisa e desenvolvimento por meio de universidades públicas e institutos de pesquisa ; como as pesquisas realizadas em outros países possibilitaram acesso mais rápido às vacinas e serviram de modelo para o que foi desenvolvido no Brasil ; como o contingenciamento do Fundo Nacional de Desenvolvimento Científico e Tecnológico ( FNDCT ) prejudicou e ainda prejudica o Brasil, entre diversos outros assuntos.
No item B, era esperado que o estudante explicasse como o fomento público ao desenvolvimento científico pode atender à justiça social e à equidade, em contextos como o da pandemia de Covid-19.
Além disso, deveria abordar a importância das universidades com a justiça social ; como os países que investem nessa área têm grandes vantagens em relação a outros países que não investem ; a distribuição de vacina atende à justiça social por possibilitar que toda a população seja vacinada, independentemente de ser rico ou pobre, por exemplo.
O espaço disponível para a resposta dos dois itens foi adequado e a maioria das respostas ocupou parte das linhas, não utilizando todo o espaço.
Algumas respostas não apresentaram a separação dos itens A e B.
A maioria dos estudantes respondeu de modo semelhante ao que constava do padrão de resposta.
Foi recorrente a exposição do problema sem desenvolvê -lo com clareza, apenas de forma superficial.
A abordagem do candidato sobre o tema tornou -se, em alguns casos, muito genérica.
As respostas consideradas boas atenderam ao que a banca esperava.
Abordaram a importância de investir em CT\&I para desenvolver o país, desdobrando os argumentos e citando como isso poderia ocorrer.
No geral, as notas estavam abaixo do esperado.
Ficou evidente que a maioria dos estudantes possui um conhecimento sobre os dois itens apresentados.
A Questão 1 visou avaliar o desenvolvimento da argumentação de um tema relevante para a sociedade brasileira vigente.
Além disso, houve problemas estruturais nos textos, pois a grande maioria dos estudantes tem uma noção do assunto, mas muitos possuem dificuldade de construir um texto com coesão e coerência, ou seja, não conseguem concatenar os argumentos para defender um ponto de vista.
No geral, em relação à Questão 1, identificou -se que a principal deficiência na formação dos estudantes está relacionada ao aprofundamento do conteúdo e à capacidade de expressar, com clareza e coesão, sua opinião a respeito de determinado assunto. 7.3.1.3 ANÁLISE DE CONTEÚDO DA QUESTÃO DISCURSIVA 2 DO COMPONENTE DE FORMAÇÃO GERAL A Tabela 7.15 apresenta dados estatísticos a respeito das notas de conteúdo da Questão Discursiva 2 do Componente de Formação geral, organizados por grande região de funcionamento dos cursos, no país.
A Tabela 7.15 mostra que o desempenho médio dos estudantes, na Questão Discursiva 2 ( Média 57,9 ), foi acima do obtido na Questão Discursiva 1 ( Média 53,9 ).
A região Centro-Oeste foi aquela cuja Média foi menor ( 55,8 ), e a maior Média foi na região Nordeste ( 58,8 ).
Quanto à variabilidade das notas, o desvio-padrão de todo o Brasil foi 26,6 \% superior ao obtido na Questão Discursiva 1 ( 25,7 ).
O maior desvio nessa questão foi obtido na região Centro-Oeste ( 27,6 ), enquanto o menor foi obtido na região Sudeste ( 26,1 ).
A mediana de todo o Brasil foi 63,5, a mesma de duas regiões.
As exceções foram as regiões Nordeste e Sul, que obtiveram Mediana de 65,5.
A nota máxima da Questão Discursiva 2 foi 100,0, obtida por pelo menos em todas as regiões.
Já a nota mínima para todas as regiões foi zero.
No Gráfico 7.8, mostra -se a distribuição das notas na Questão Discursiva 2 do Componente de Formação geral.
O Gráfico 7.8 mostras que pouco mais de 10 \% dos participantes deixaram em branco a resposta à questão discursiva 2.
A moda dessa distribuição ocorre no intervalo ( 70 ; 80 ] com uma frequência maior que 20 \% dos participantes. 7.3.1.4 COMENTÁRIOS SOBRE A CORREÇÃO DE CONTEÚDO DAS RESPOSTAS À QUESTÃO DISCURSIVA 2 A Questão Discursiva 2 abordou um tema importante, com o objetivo de analisar conhecimentos gerais do aluno.
O tema contemplou o conhecimento geral do patriarcado estrutural ( dominação masculina ), o qual é recorrente nos noticiários, e um infortúnio na vida dos brasileiros, pois trata -se de uma questão de violência contra as mulheres.
É recorrente assistir ou ler notícias sobre o domínio masculino na sociedade brasileira.
Questões como a participação da mulher na política ; a paridade de salários entre profissionais de sexos diferentes ; a desigualdade social e o tratamento hostil que é dado às mulheres, em diversas situações, tanto profissionais quanto familiares, devem ser de conhecimento geral do estudante.
Não se objetiva tratar questões de ordem específica para a solução deste problema, mas, apenas, discorrer e argumentar sobre ele.
Portanto, a questão foi elaborada de maneira oportuna e adequada para o universo do estudante universitário brasileiro.
A Questão está dividida em dois itens : A e B.
No item A, esperava -se que o estudante abordasse o patriarcado – dominação masculina que gera a divisão sexual do trabalho, a desigualdade econômica e política, entre homens e mulheres, a sub-representação das mulheres na política e em cargos de poder.
Esperava -se que o estudante ressalta se, ainda, questões culturais, como a educação sexista e a mídia, que coisificam a imagem das mulheres e naturalizam a desigualdade entre os sexos, de forma a considerar os homens como dominantes, enquanto as mulheres são historicamente tidas como submissas e incapazes.
No item B, era esperado que o estudante apresentasse a importância do papel do Estado na implementação de políticas públicas voltadas à proteção da mulher, tais como na área da saúde, na educação não sexista, na assistência social, na segurança pública – delegacias femininas, defensoria pública especializada, casa de abrigo, juizado especial etc., bem como ONGs que atuam no combate à violência contra a mulher e ao feminicídio.
O espaço disponível para a resposta dos dois itens foi adequado e grande parte das respostas ocupou parte das linhas, não contemplando o espaço completo.
Com relação à linguagem utilizada pelos estudantes, as respostas estavam dentro do esperado.
Alguns estudantes utilizavam termos do jargão do cotidiano.
No item B, o estudante deveria propor duas ações do Estado para o enfrentamento da violência contra as mulheres.
Os alunos deveriam apresentar propostas e desenvolvê -las.
Muitos estudantes responderam apenas por meio da apresentação de ações do Estado, mas sem desenvolvê -las ou argumenta -las.
Em geral, os alunos apenas expuseram ações – algumas vezes, uma, duas e, até mesmo, três ações –, mas sem fazer relações entre elas ou sem desenvolvê -las.
As respostas dos estudantes demonstraram que a maioria apresenta um conhecimento geral a respeito do patriarcado estrutural e da trágica desigualdade que esta impõe dos homens sobre as mulheres, bem como de possíveis ações do Estado para a proteção delas.
A Questão visou avaliar o desenvolvimento da argumentação de um tema relevante para a sociedade brasileira vigente. 7.3.1.5 CONSIDERAÇÕES FINAIS SOBRE A CORREÇÃO DE CONTEÚDO DAS QUESTÕES DISCURSIVAS DO COMPONENTE DE FORMAÇÃO GERAL As questões discursivas da área de Formação Geral abordaram temas atuais e proporcionaram a avaliação de competências importantes dentre as previstas pelas suas Diretrizes Curriculares.
A análise das respostas revelou, principalmente, dificuldade dos estudantes em relação à seleção de argumentos sólidos para discorrer sobre o tema – que pode ser fruto de desinteresse ou de desconhecimento mais profundo do contexto.
Esses fatos parecem ter sido os principais responsáveis pelo resultado aquém do esperado em termos de avaliação quantitativa.
Mesmo os estudantes que responderam em conformidade com o padrão demonstraram um universo de argumentação muito limitado, pobre em vocabulário, com o recurso frequente de procurar nos textos introdutórios os elementos para a construção de suas frases.
Encadear e relacionar argumentos foi uma deficiência percebida.
As razões eram apresentadas de forma confusa, sem fluência e sem constituir unidade, utilizando termos genéricos e vagos. 7.3.1.6 ANÁLISE DE LÍNGUA PORTUGUESA DAS QUESTÕES DISCURSIVAS DO COMPONENTE DE FORMAÇÃO GERAL Os dados de Direito obtidos a partir das respostas às questões discursivas do Componente de Formação geral, no que tange à Língua Portuguesa, são apresentados na Tabela 7.16.
Os dados da Tabela 7.16 mostram que, em relação ao quesito Língua Portuguesa, os estudantes, de todo o Brasil, obtiveram Média 63,2.
A maior média em relação à Língua Portuguesa foi obtida na região Sul ( 63,8 ), e a menor, na região Centro-Oeste ( 61,5 ).
Quanto à variabilidade das notas, o desvio-padrão de todo o Brasil foi 24,5.
O menor desvio-padrão foi obtido na região Sudeste ( 24,1 ), e o maior desvio-padrão foi obtido na região Centro-Oeste ( 25,8 ).
A mediana das notas de Língua Portuguesa foi 67,5 para o Brasil, como um todo, a mesma obtida em todas as regiões.
A nota máxima para todo o Brasil foi 100, com, pelo menos, um estudante tirando essa nota em todas as regiões.
A nota mínima foi zero em todas as regiões do país.
No Gráfico 7.9, mostra -se a distribuição das notas de Língua Portuguesa do Componente de Formação geral.
O Gráfico 7.9 mostra que mais de 7 \% dos participantes não pontuaram em língua portuguesa nas questões discursivas do componente de formação geral.
A moda dessa distribuição ocorre no intervalo ( 70 ; 80 ] com uma frequência maior que 20 \% dos participantes. 7.3.1.7 COMENTÁRIOS SOBRE A CORREÇÃO DE LÍNGUA PORTUGUESA DAS RESPOSTAS ÀS QUESTÕES DISCURSIVAS Os comentários feitos neste item a respeito da correção de Língua Portuguesa das respostas emitidas para as questões discursivas do Componente de Formação Geral estão agrupados em três eixos : avaliação de desempenho linguístico ; padrão de resposta utilizado na avaliação e análise das respostas dadas pelos estudantes. a ) A avaliação de desempenho linguístico De acordo com a Portaria n.º 298 do MEC/INEP, de 30 de junho de 2022, o Enade deve avaliar se o estudante desenvolveu, no processo de formação, dentre outras, competências para [ … ] V – Compreender as linguagens e suas respectivas variações ; VI – Ler, produzir e interpretar textos com clareza e com coerência ; VII – Identificar, analisar e interpretar diferentes representações verbais, não verbais, gráficas e numéricas de fenômenos diversos ou de um mesmo significado ; VIII – Formular e articular argumentos e contra – argumentos consistentes em situações sociocomunicativas ; IX – Identificar, compreender e analisar situações-problema utilizando pensamento holístico e sistêmico ao se abordar a complexidade da realidade. ( INEP, 2022d ) Essas competências são essenciais à formação de nível superior, porque dizem respeito ao domínio da linguagem, tanto no processo de leitura quanto na articulação e na organização de ideias para a produção textual.
Dessa forma, as duas questões discursivas de formação geral oferecem oportunidade para avaliar se o participante consegue desenvolver um tema, recorrendo a argumentos que explicitem sua opinião e elaborando propostas de intervenção social adequadas a um exercício consciente da cidadania.
É, portanto, necessário identificar, nas respostas discursivas de Formação geral, as diferentes habilidades necessárias ao bom desempenho linguístico na situação de formalidade exigida : a organização das ideias, tendo em vista a construção de um texto coerente e estruturado ; o domínio das convenções ortográficas e das convenções morfossintáticas que constituem o arcabouço da gramática da norma-padrão da Língua Portuguesa ; e o domínio de um vocabulário amplo e diversificado, que permita uma comunicação clara e eficiente.
A Questão Discursiva 1 desenvolveu o tema de investimentos públicos em educação, ciência, tecnologia e inovação, que são estratégicos e contribuem para o desenvolvimento científico de um país, podendo abordar essa questão tanto no Brasil quanto em outros países.
O enunciado estava claro ao encaminhar a solicitação de resposta, evidenciando a exigência de elaboração de um texto que discutisse e explicasse por que os investimentos públicos em educação, ciência, tecnologia e inovação são estratégicos e contribuem para o desenvolvimento científico de um país e como o fomento público ao desenvolvimento científico pode atender à justiça social e à equidade, em contextos como o da pandemia de Covid-19.
Os três textos motivadores favoreceram a formulação e a articulação de argumentos e contra-argumentos consistentes, em defesa do ponto de vista a ser defendido pelo participante.
A Questão Discursiva 2 desenvolveu o conceito de patriarcado, tomando como base um texto de Biroli.
Com base nesse texto, o comando solicitou do participante :"discorra sobre a relação entre o patriarcado estrutural e a desigualdade entre homens e mulheres ” e “ proponha duas ações do Estado para o enfrentamento dos diferentes tipos de violência contra a mulher ”.
O comando “ discorra ” demandava a elaboração de um texto de base dissertativa, com desenvolvimento de argumentação sobre o tema.
Em relação à segunda parte, o comando exigia a elaboração de duas propostas de ações, que deveriam conter informações sobre o agente da intervenção ( Estado ) para o enfrentamento da situação apresentada.
As duas questões, portanto, assemelhavam -se, pois, solicitavam um texto dissertativo ao apresentarem comandos para “ discutir ” e “ discorrer ”.
A Questão Discursiva 2, entretanto, solicitava que o aluno apresentasse propostas de ações para o enfrentamento do problema destacado.
Esses dois estilos textuais ofereceram oportunidade de avaliar se os participantes estavam habilitados a desenvolverem textos com esses dois distintos objetivos comunicativos, o que seria esperado ao final do curso de graduação e necessário à sua vida profissional.
A análise do desempenho linguístico esperado de um estudante que está concluindo seu curso superior apoiou -se no princípio de que ele deveria utilizar conhecimentos sobre os temas e estruturar seus textos de acordo com as características do registro formal adequado à situação comunicativa – avaliação de conhecimentos ao final de um curso superior.
Essa configuração determina exigências quanto à adequação da seleção vocabular, ao desenvolvimento do conteúdo, à estruturação sintática dos períodos, à organização lógica das ideias, à utilização de procedimentos de encadeamento textual e de referenciação, à obediência às exigências morfossintáticas próprias da modalidade escrita da norma-padrão, ao respeito às regras ortográficas de acentuação gráfica e de grafia das palavras. a ) O padrão de resposta utilizado na avaliação Em relação a este eixo de análise, foram considerados, no padrão de resposta, as seguintes competências : ( 1 ) estruturação textual ; ( 2 ) respeito às convenções ortográficas da norma-padrão da Língua Portuguesa ; ( 3 ) domínio dos diferentes aspectos morfossintáticos próprios da modalidade escrita formal da norma-padrão da Língua Portuguesa ; e ( 4 ) seleção vocabular adequada à modalidade ( 1 ) Estruturação textual condizente com o gênero textual solicitado e o modo de organização textual adequado ao gênero ( argumentativo ou expositivo ).
Esta competência envolve : - estruturação sintática condizente com o padrão da modalidade escrita formal da Língua Portuguesa, de modo a garantir a clareza e a coerência necessárias ; - distribuição do conteúdo do texto em parágrafos, de modo a garantir a sua organização temática ; - distribuição das ideias em períodos que mantenham uma estruturação sintática condizente com a modalidade formal da língua ( evitando -se frases fragmentadas, sequência justaposta de ideias sem encaixamentos sintáticos e presença de oração subordinada sem oração principal que reflitam hábitos da oralidade e comprometam a estrutura lógico-gramatical do texto ) ; - utilização de operadores discursivos que contribuam para a progressão temática do texto, estabelecendo relações lógicas entre as ideias apresentadas, tanto do ponto de vista intrafrasal, como interfrasal : preposições, conjunções, pronomes relativos, advérbios e locuções adverbiais de tempo e lugar ; - utilização de procedimentos de referenciação lexical e pronominal que permitam a retomada de referentes textuais, como pronomes, advérbios e artigos ; além de elementos lexicais, como sinônimos, hiperônimos, hipônimos, expressões metafóricas e expressões resumitivas ) ; - emprego adequado dos pronomes relativos “ cujo(a) ” e “ onde ”, responsáveis por encaixamentos sintáticos no período composto ; - utilização adequada de sinais de pontuação que contribuam para a organização lógica da frase e do texto ; - inteligibilidade relacionada ao atendimento das exigências de estruturação textual.
A expectativa, portanto, era de que o participante recorresse a procedimentos linguístico–discursivos para organizar seu texto, permitindo o encadeamento lógico entre suas partes de forma a garantir a progressão e a coerência textuais, com estruturação que não reproduzisse hábitos da oralidade e sim atendesse à norma-padrão da Língua Portuguesa. ( 2 ) Respeito às convenções ortográficas da norma-padrão da Língua Portuguesa.
Esta competência envolve o respeito às convenções estabelecidas pela legislação em vigor e consubstanciadas no Vocabulário Ortográfico da Língua Portuguesa, editado pela Academia Brasileira de Letras e avaliam -se : - domínio das regras de acentuação gráfica ; - domínio da grafia padrão das palavras ( com ausência de abreviaturas próprias da linguagem da internet ).
Assim, esperava -se que o participante : grafasse corretamente as palavras ; respeitasse as regras de acentuação gráfica ; empregasse maiúsculas em início de frase, em nomes próprios de pessoas, lugares ou instituições ; grafasse as siglas com letras maiúsculas ou apenas com a primeira letra maiúscula quando formarem uma palavra ; evitasse abreviações como p/, vc, tb, pq, tá, né, usadas muitas vezes em escrita informal e na internet ; obedecesse às regras de separação de sílabas no final da linha ; atendesse aos padrões de segmentação vocabular da norma-padrão. ( 3 ) Domínio dos diferentes aspectos morfossintáticos próprios da modalidade escrita formal da norma-padrão da Língua Portuguesa.
Esta competência envolve o domínio das convenções próprias da norma padrão da Língua Portuguesa relativas aos seguintes aspectos : concordância nominal ; concordância verbal ; regência nominal ; regência verbal ; flexão nominal ; flexão verbal ; correlação entre os tempos verbais ; colocação pronominal. ( 4 ) Assim, esperava -se que o participante : - flexionasse o verbo para estabelecer concordância de número com o sujeito da frase ; - flexionasse o artigo, o adjetivo e o pronome para concordar em número e em gênero com o substantivo a que cada um se refere ; - observasse a regência nominal e a verbal, utilizando a preposição adequada depois de um substantivo, um verbo ou um adjetivo ; - empregasse adequadamente o acento grave indicador da crase que expresse a contração entre uma preposição e um artigo, de acordo com as exigências de regência nominal ou verbal ; - obedecesse às regras de colocação pronominal ( próclise e ênclise ), distintas dos hábitos da oralidade ou da escrita informal ; - flexionasse adequadamente verbos, substantivos, adjetivos e pronomes, no que diz respeito à expressão das categorias gramaticais ; - flexionasse os verbos para expressar a correlação de modo e de tempo nas estruturas subordinadas.
Com base nesses critérios, alguns desvios ortográficos foram considerados como desvios de caráter morfossintático porque interferem na classificação dos vocábulos de base morfológica e funcional : - eliminação da marca de infinitivo ( – r – ) e substituição por acento agudo ou ausência total de marca do infinitivo, causando ambiguidade entre distintas classes de palavras ; - confusão entre as terminações “ ão ” e “ am ”, utilizadas para marcar diferentes tempos verbais ; - emprego indevido da forma verbal de"haver"\&lt ; “ há”\&gt ; no lugar da preposição ou do artigo \&lt ; “ a”\&gt ;, e vice-versa ; - emprego de hífen para separar pronome átono – tanto uso indevido quanto omissão ( ex : “ esperasse ”, em lugar de “ espera -se ” ; “ falar–mos ”, no lugar de “ falarmos ” ) ; - emprego de “ esta ” no lugar de “ está ”, provocando ambiguidade devido à confluência gráfica entre o pronome demonstrativo e a forma do verbo"estar", pela ausência da acentuação gráfica ; - emprego da conjunção aditiva"e"no lugar da forma verbal “ é ”, provocando ambiguidade devido à confluência gráfica por ausência da acentuação gráfica ; - emprego do advérbio “ mais ” no lugar da conjunção adversativa “ mas ”, reproduzindo confluência fônica da oralidade ; - registro da terceira pessoa do plural dos verbos “ ter ” e “ vir ” sem o acento circunflexo, causando desvio de concordância verbal ; - emprego inadequado do relativo “ onde ” deslocado do seu sentido locativo, reproduzindo tendência da oralidade ; - emprego do pronome relativo “ que ” sem a necessária preposição ( desvio de regência ) ; - ausência do acento grave obrigatório para marcar a crase. ( 5 ) Seleção vocabular adequada à modalidade escrita formal da Língua Portuguesa, exigida pela situação comunicativa.
Esta competência envolve : - precisão na seleção/utilização do vocabulário relacionado à temática solicitada pela questão ; - diversificação e adequação vocabular ; - coerência na combinação vocabular ; - ausência de marcas da oralidade, como termos de sentido muito genérico ( “ coisa ”, “ negócio ”, “ você ” ) e termos de registros mais informais ( como gírias, jargões, frases feitas, ditados populares, termos regionais ).
Assim, esperava -se que o participante : respeitasse a adequação vocabular, não empregando gírias ou expressões coloquiais ; apresentasse repertório vocabular diversificado e preciso ; evitasse repetição desnecessária de palavras ; utilizasse um vocabulário mais formal, como solicitado por um texto dissertativo.
Essas quatro competências foram reunidas em três itens de avaliação, atendendo aos aspectos explicitados anteriormente : I – Aspectos Ortográficos ( 30 \% ) : domínio das convenções ortográficas : grafia de vogais e consoantes, uso de maiúsculas e minúsculas, emprego do hífen e da acentuação gráfica, segmentação de palavras ; II – Aspectos Textuais ( 20 \% ) : domínio dos procedimentos de estruturação textual : organização interna dos períodos, emprego de conectores para a articulação lógica entre os períodos e entre os parágrafos, emprego de marcas de referenciação lexical e pronominal ; utilização dos sinais de pontuação que contribuem para a organização lógica da frase ; III – Aspectos Morfossintáticos e Vocabulares ( 50 \% ) : domínio das regras de caráter morfossintático estabelecidas como modelares do ponto de vista da modalidade escrita formal da norma-padrão da Língua Portuguesa : concordância nominal e verbal, regência nominal e verbal, colocação pronominal, flexão nominal e verbal, correlação entre tempos e modos verbais, ausência de marcas de oralidade.
Outra competência combinada a esta na análise das respostas é a seleção vocabular adequada à modalidade escrita formal da Língua Portuguesa, tendo em vista a intersecção entre as duas do ponto de vista das exigências do registro formal da modalidade escrita da norma-padrão. b ) Análise das respostas dadas pelos estudantes A Questão 1 propiciou a oportunidade de maior desenvolvimento do tema solicitado, gerando consequências na estruturação textual das respostas, já que os textos foram mais expositivos, mais elaborados e fluentes, inclusive, um dos textos era uma tabela.
Já a Questão 2 teve um comportamento distinto, devido ao seu caráter mais opinativo, o texto foi mais curto e o vocabulário mais limitado.
Observaram -se, nas duas questões, alguns casos de comprometimento de autoria, com textos muito próximos a paráfrases dos textos motivadores e sem marcas de repertório sociocultural produtivo.
Apresentam -se, a seguir, algumas \textbf{observações} sobre o desempenho dos participantes : I – Aspectos Ortográficos : o desempenho dos participantes revelou uma diferença muito grande nos dois aspectos analisados, como já foi observado nos relatórios dos anos anteriores : baixo índice de desvios da grafia padrão e grande índice de desvios de acentuação.
Em vários casos, ocorre ausência quase completa de acentuação gráfica.
Vale ressaltar que muitos participantes apresentaram uma caligrafia quase ininteligível.
Ainda nesse aspecto, ocorreu a omissão inadequada de cedilha, de til e do traço do T em algumas respostas.
Houve, inclusive, respostas com ausência total de cedilha ou til.
Os casos mais sistemáticos de eliminação do acento indicador da sílaba tônica foram : palavras proparoxítonas ( “ politica ”, “ pratica ”, “ publica ”,"policia") ; palavras paroxítonas terminadas em ditongo crescente ( “ necessario ”,"historia","propria", “ violencia ” ) ; palavras oxítonas ( “ tambem ”,"populacao", “ pais ” ).
Por outro lado, destaca -se o uso indevido do acento gráfico em palavras como “ investímento ” e “ pátriarcado ”, por exemplo.
Quanto aos desvios de grafia, observam -se os seguintes exemplos : emprego inadequado de maiúscula : “ Leis ” ; ou de minúscula :"constituição"“ Ministério publico ” ; confusão entre"ss","sc","c"e"ç": “ nessessaria", “ agreção ”, “ conciencia ”, “ comecam ”, “ cocientização ”, “ pacos ” ; eliminação de grafemas : “ mostando ”,"deveriamo"; - desvios de segmentação : “ des de", “ des igualdade ”, “ pre conceito"; troca de grafema : “ órgam ” ( por “ órgão ” ), “ escravido ” ( por “ escravidão ” ).
II – Aspectos Textuais : a estruturação textual foi a competência que apresentou o maior número de desvios, sendo responsável, em alguns textos, pela falta de clareza e pela desorganização.
No conjunto dos textos analisados, observaram -se diferentes níveis de domínio dos recursos de estruturação e organização textual.
No entanto, os participantes podem ser agrupados em três faixas descritas a seguir : - os que possuem grande domínio dos recursos de estruturação textual, com produção de textos que atendem às exigências de clareza e coerência, com progressão temática bem direcionada e utilização produtiva de operadores argumentativos ; - os que apresentam um domínio regular dos recursos de estruturação textual, com desvios na organização do texto e ausência de paragrafação, além de uso insuficiente de operadores argumentativos e processos de encaixamento sintático ; - os que apresentam baixo domínio dos recursos de estruturação textual, com muitos desvios que retomam hábitos da oralidade e que comprometem a coerência do texto : ocorrência de sequência justaposta de ideias sem encaixamentos sintáticos ; redução de estruturas subordinadas, ao lado do aumento na frequência de estruturas coordenadas e absolutas ; redução no uso de conectores para expressar relações lógicas essenciais à construção do texto ; emprego equivocado de operadores que não estabelecem relações lógicas coerentes entre ideias do texto ; fragmentação frasal que compromete a estrutura lógico–gramatical.
Uma estrutura que vem se ampliando na escrita dos participantes, refletindo hábitos da oralidade, é a duplicação de um dos elementos da frase, ou seja, a retomada anafórica de um elemento topicalizado, como por exemplo :"Elas precisam saber o verdadeiro papel DELAS na sociedade “, A desigualdade ELA vem de uma cultura machista... ” ; “ Homem e mulher, ELES são iguais ”.
Também em relação ao emprego dos sinais de pontuação, observaram -se distintos tipos de desempenho : participantes que dominam a funcionalidade da pontuação para a construção de textos fluentes e coesos, por um lado, e um grande número de participantes que apresentam precariedade em sua utilização, chegando ao ponto de elaborar parágrafos sem qualquer marca interna de pontuação para separar os períodos.
São os seguintes os tipos de desvios encontrados : - emprego de vírgula para separar o sujeito e o predicado : “ A mulher, sofre violência ” ( no lugar de “ A mulher sofre violência ” ) ;"O correto seria, ter um canal de denúncia ” ( no lugar de “ o correto seria ter um canal de denúncia ” ) ; - ocorrência de apenas uma das vírgulas para separar uma palavra, uma expressão ou uma oração encaixada : “ O homem, por sua vez não aceita que uma mulher seja sua chefa ” ( no lugar de “ O homem, por sua vez, não aceita que uma mulher seja sua chefa ” ) ; - emprego de ponto no lugar de vírgula para separar orações de um mesmo período : “ A todo custo as mulheres são atacadas na política.
Sexualmente também. ” ( no lugar de “ A todo custo as mulheres são atacadas na política, sexualmente também. ” ) ; - ausência de vírgula ou conjunção para separar elementos de uma enumeração :"... conquistas como voto equidade salarial direito de estudar e reconhecimento..."( no lugar de""... conquistas como voto, equidade salarial, direito de estudar e reconhecimento...") ; - ausência de vírgula para separar oração adjetiva explicativa ou utilização inadequada para separar oração adjetiva restritiva : “ As mulheres, que são iguais aos homens devem ser igualmente protegidas por lei."( no lugar de “ As mulheres, que são iguais aos homens, devem ser igualmente protegidas por lei.") ; - emprego indevido do sinal de dois pontos : “ Mulheres são julgadas : pelas simples roupas que usam. ” ( no lugar de “ Mulheres são julgadas pelas simples roupas que usam. ” ) ; - ausência de pronome no processo de referenciação : “ julga a mulher ou até agride ” ( no lugar de “ julga a mulher ou até a agride ” ” ) ; - emprego de vírgula em lugar do ponto – final : “ O patriarcado estrutural é formado por elementos históricos e religiosos, há pouquíssimo tempo, as mulheres passaram a votar. ” ( no lugar de “ O patriarcado estrutural é formado por elementos históricos e religiosos.
Há pouquíssimo tempo, as mulheres passaram a votar. ” ).
III – Aspectos Morfossintáticos e Vocabulares : o domínio das convenções morfossintáticas apresentou resultados que evidenciam, claramente, alguns aspectos problemáticos no desempenho dos participantes.
Assim como na análise dos desvios textuais, pode -se observar uma distinção entre três grupos de participantes : o grupo que revela domínio das convenções morfossintáticas ; o grupo que apresenta um domínio regular, com desvios pontuais ; e o grupo que revela domínio precário dessas convenções.
Pode -se destacar a regência, tanto verbal quanto nominal, como o aspecto com maior número de desvios : a ) Regência nominal e verbal : Na maioria dos textos, destacam -se os desvios relacionados à noção de"regência", ou seja, à relação de subordinação que se estabelece entre as palavras por meio, ou não, de preposições, refletindo uma tendência na modalidade escrita da língua.
O desvio mais frequente foi a falta do sinal indicativo da crase, em contextos em que seria obrigatória, ou, ao contrário, um uso indevido em contextos em que o princípio não se aplica.
Os dois tipos de desvios revelam que o usuário não tem consciência de que, sob a forma do termo “ a ”, se pressupõe a presença de uma preposição “ a ”, exigida pela regência do termo anterior.
Exemplos : “ diferentes tipos de agressão a mulher ” ( no lugar de “ diferentes tipos de agressão à mulher ” ) ; “ Desenvolverem práticas econômicas e sociais as novas gerações"( no lugar de “ Desenvolverem práticas econômicas e sociais às novas gerações") ; “ Além de proteger à mulher... ” ( no lugar de “ Além de proteger a mulher... ” ).
Outro desvio relacionado à regência verbal e à nominal, encontrado frequentemente nas questões, é a ausência de preposição antes de pronome relativo, refletindo um processo generalizado na modalidade oral da língua, em situações de registro informal.
Exemplo :"outra forma é oferecer apoio psicológico quem sofreu qualquer tipo de abuso..."( no lugar de :"outra forma é oferecer apoio psicológico a quem sofreu qualquer tipo de abuso...") b ) Concordância nominal e verbal : Os desvios relativos à aplicação das regras de concordância verbal e nominal correspondem a fatos que têm sido amplamente estudados na modalidade oral da Língua Portuguesa.
Quanto à concordância de número, observou -se ausência de marca ( com sujeito anteposto ou posposto ), como por exemplo a ausência de acento circunflexo na forma plural do presente do indicativo do verbo “ ter ”, considerada como um desvio na concordância verbal e não na acentuação gráfica.
Exemplo :"as mulheres ainda tem como obrigação"( no lugar de"as mulheres têm como obrigação") ; “ as leis em vigor tem que ser cumpridas ” ( no lugar de “ as leis em vigor têm que ser cumpridas ” ).
Outro desvio na aplicação da regra da concordância de número foi o aparecimento da marca de plural em verbos ou adjetivos relacionados a núcleos substantivos no singular, devido à presença de modificadores no plural, evidenciando um processo de hipercorreção.
Esse processo tem se ampliado nas situações de oralidade, caracterizando uma tendência de variação da Língua Portuguesa.
Exemplos :"Mudança nas leis que aumentem a pena do agressor"( no lugar de"Mudança nas leis que aumente a pena do agressor") ;"o aumento dos feminicídios devem ser combatidos ” ( no lugar de"o aumento dos feminicídios deve ser combatidos ” ). tipos de desvios revelam que o usuário não tem consciência de que, sob a forma do termo “ a ”, se pressupõe a presença de uma preposição “ a ”, exigida pela regência do termo anterior.
Exemplos : “ diferentes tipos de agressão a mulher ” ( no lugar de “ diferentes tipos de agressão à mulher ” ) ; “ Desenvolverem práticas econômicas e sociais as novas gerações"( no lugar de “ Desenvolverem práticas econômicas e sociais às novas gerações") ; “ Além de proteger à mulher... ” ( no lugar de “ Além de proteger a mulher... ” ).
Outro desvio relacionado à regência verbal e à nominal, encontrado frequentemente nas questões, é a ausência de preposição antes de pronome relativo, refletindo um processo generalizado na modalidade oral da língua, em situações de registro informal.
Exemplo :"outra forma é oferecer apoio psicológico quem sofreu qualquer tipo de abuso..."( no lugar de :"outra forma é oferecer apoio psicológico a quem sofreu qualquer tipo de abuso...") b ) Concordância nominal e verbal : Os desvios relativos à aplicação das regras de concordância verbal e nominal correspondem a fatos que têm sido amplamente estudados na modalidade oral da Língua Portuguesa.
Quanto à concordância de número, observou -se ausência de marca ( com sujeito anteposto ou posposto ), como por exemplo a ausência de acento circunflexo na forma plural do presente do indicativo do verbo “ ter ”, considerada como um desvio na concordância verbal e não na acentuação gráfica.
Exemplo :"as mulheres ainda tem como obrigação"( no lugar de"as mulheres têm como obrigação") ; “ as leis em vigor tem que ser cumpridas ” ( no lugar de “ as leis em vigor têm que ser cumpridas ” ).
Outro desvio na aplicação da regra da concordância de número foi o aparecimento da marca de plural em verbos ou adjetivos relacionados a núcleos substantivos no singular, devido à presença de modificadores no plural, evidenciando um processo de hipercorreção.
Esse processo tem se ampliado nas situações de oralidade, caracterizando uma tendência de variação da Língua Portuguesa.
Exemplos :"Mudança nas leis que aumentem a pena do agressor"( no lugar de"Mudança nas leis que aumente a pena do agressor") ;"o aumento dos feminicídios devem ser combatidos ” ( no lugar de"o aumento dos feminicídios deve ser combatidos ” ).
Quanto à concordância de gênero, vários casos foram observados, normalmente no âmbito de sintagmas nominais longos, em que o adjetivo está afastado do substantivo, conforme apontam as pesquisas sobre a modalidade oral.
Exemplo : “ E as leis sobre as mulheres que sofrem violência ser mais rigorosos ” ( no lugar de “ E as leis sobre as mulheres que sofrem violência ser mais rigorosas ” ) ; Registro da forma verbal “ é ” sem acento, provocando ambiguidade na leitura pela confluência com a conjunção aditiva"e".
Exemplo : “ A desvantagem entre salários, poderes e entre a violência física e moral e crescente por demais... ” ( no lugar de “ A desvantagem entre salários, poderes e entre a violência física e moral é crescente por demais ”. c ) Emprego do pronome relativo “ onde ” como relativo universal, típica marca da oralidade, em que esse pronome não é mais utilizado apenas em sua função locativa"lugar onde se realiza a ação".
Exemplo :"Existe a desigualdade de gênero, onde homem ganha mais do que as mulheres"; “ há falta de informação e conhecimento onde ” ; “ O patriarcado está estruturado em nossa sociedade há séculos, onde o sexo masculino... ” d ) Falta de artigo definido antes de substantivo, comprometendo o paralelismo em enumerações : Exemplo :"seja em oportunidade de emprego, salário"( no lugar de"seja em oportunidade de emprego, de salário"). e ) Desvios relacionados à flexão verbal : Exemplos : emprego indevido do tempo ou do modo verbal : “ promovido ” por “ promovendo ” ; eliminação de marca de infinitivo : “ colocar ” por “ coloca ”, “ ta ” por “ estar ”, “ está ” por “ estar ”, “ esta ” por “ estar ” ; “ estás ” por “ estas ” ; utilização de construção de “ estar ” + gerúndio como marca de oralidade : “ a sociedade cada vez mais tá evoluindo ” ; segmentação indevida da desinência de número e pessoa de verbos : “ denunciar – mos ”.
Assim como nas demais competências avaliadas, há diferenças entre os participantes em relação à seleção e ao emprego do vocabulário.
Foram encontrados textos com alto padrão de domínio vocabular, chegando até ao emprego de vocabulário técnico das áreas do saber abordadas nas questões.
Por outro lado, na maioria dos textos, foram observados vários tipos de inadequação na seleção e na utilização do vocabulário, sendo que o principal deles foi a excessiva repetição de certas palavras, revelando limitação de repertório vocabular.
Além disso, observou -se a ocorrência de expressões da oralidade, revelando falta de domínio de vocabulário mais abstrato e de maior complexidade, essencial ao desenvolvimento do texto de base dissertativa.
Exemplos : seleção vocabular indevida : “ Em pleno século 21 é de em plena importância... ” ; emprego inadequado do verbo “ ter ” no lugar de “ haver ” : “ Durante a história, teve muito machismo ” ; imprecisão vocabular : “ Mas a cabeça, o corpo, há mesma igualdade para desenvolvermos igualmente as coisas ”, “ Em alguns anos atrás, as coisas eram mais complicadas nas questões de desigualdade no trabalho. ”.
A imprecisão vocabular, como o uso da palavra “ coisa ” em vários textos, demonstra o domínio precário da norma-padrão da Língua Portuguesa, com presença constante de traços de oralidade. 7.3.1.8 CONSIDERAÇÕES FINAIS SOBRE A CORREÇÃO DE LÍNGUA PORTUGUESA DAS RESPOSTAS ÀS QUESTÕES DISCURSIVAS DO COMPONENTE DE FORMAÇÃO GERAL O conjunto de textos analisados no Enade 2022 permite traçar um perfil bem detalhado do desempenho linguístico do estudante brasileiro de graduação.
Conforme ficou comprovado nos comentários dos exemplos de respostas apresentados neste relatório, uma parcela significativa dos participantes revela lacunas importantes em sua formação, principalmente no que diz respeito à estruturação e à organização de textos.
Como esta é uma competência essencial à produção escrita, que deve ser construída ao longo dos anos de escolaridade ( do ensino fundamental à educação superior ), sua recuperação é trabalhosa e depende de um esforço intenso de exposição a textos e de reflexão sobre os requisitos que garantem a clareza e o encadeamento das ideias.
A competência relacionada ao domínio do vocabulário, que deve ser diversificado e apropriado ao padrão formal da língua, também apresentou um desempenho insatisfatório pela maioria dos participantes, que se limitam ao senso comum, com vocabulário genérico e vago, o que é inadequado para atender ao perfil de um graduando.
As outras duas competências avaliadas ( convenções ortográficas, convenções morfossintáticas ) apresentam desvios pontuais e são de mais fácil recuperação, com exceção dos desvios morfossintáticos que já estão incorporados na modalidade oral informal da Língua Portuguesa, configurando processos de variação linguística e, talvez, antecipando futuras mudanças no padrão da modalidade escrita.
Com base na avaliação do desempenho linguístico aqui relatada, conclui -se que é preciso aperfeiçoar a formação superior, oportunizando aos estudantes mais acesso às ferramentas necessárias para que sejam bem-sucedidos nas diferentes situações de sua vida cidadã – profissional, cultural e política.
Essa é uma contribuição efetiva que os resultados do Enade podem oferecer às IES do país.
Para finalizar, vale observar que as duas questões discursivas atenderam às exigências de avaliação, tanto do domínio discursivo dos estudantes, quanto do seu conhecimento sobre temas relacionados ao conhecimento acadêmico, o que pode ser extremamente útil para que os cursos superiores possam contribuir na formação de seus graduandos. 7.3.2 COMPONENTE ESPECÍFICO A Tabela 7.17 apresenta dados sobre questões discursivas do Componente de Conhecimento Específico.
A partir dos dados registrados na Tabela 17.7, observa -se que a Média foi mais baixa do que para as questões discursivas do Componente de Formação geral.
Enquanto, no Componente de Formação geral, a média para estudantes de Direito de todo o Brasil foi 56 ; na parte de Componente Específico, a média foi 19,8.
A maior média deste componente foi obtida pelos estudantes da região Sul ( 20,3 ), e a menor, pelos da região Centro-Oeste ( 18,4 ).
Quanto à variabilidade das notas, o desvio-padrão de todo o Brasil foi 13,6.
O maior desvio-padrão foi encontrado na região Nordeste ( 14 ), e o menor, na região Sul ( 13,3 ).
No Gráfico 7.10, representa -se a distribuição das notas nas questões discursivas no Componente de Específico.
À luz dos dados do Gráfico 7.10, verifica -se que a moda da distribuição das notas nas questões discursivas no Componente Específico ocorre no intervalo ( 0 ; 10 ], com 25 \% do total de participantes. 7.3.2.1 ANÁLISE DA QUESTÃO DISCURSIVA 3 DO COMPONENTE ESPECÍFICO Os resultados da Questão Discursiva 3 do Componente Específico são descritos na Tabela 7.18 e comentados na sequência.
Como mostrado na Tabela 7.18, na Questão 3, a média dos estudantes de todo o Brasil foi 21.
O desempenho dos estudantes nessa questão foi superior aos das questões 2 e 5.
A menor Média na questão, em evidência, foi obtida pelos estudantes da região Centro-Oeste ( 19,6 ), enquanto a maior Média foi obtida nas regiões Sul e Sudeste ( 21,5 ).
Quanto à variabilidade das notas, o desvio-padrão de todo o Brasil foi 17,7.
O maior desvio-padrão foi obtido na região Nordeste ( 18,1 ), enquanto o menor foi obtido na região Sul ( 17,4 ).
A nota máxima, 100,0 pontos, foi alcançada por, pelo menos, um estudante em cada região.
A mediana do Brasil foi 17,5, para as regiões Centro-Oeste, Nordeste e Norte ; nas demais regiões foi de 20.
A nota mínima zero foi a mesma em todas as regiões do Brasil.
Destaca -se, com base nos dados do Gráfico 7.11, que mais de 20,0 \% dos estudantes deixaram a resposta à questão em branco.
A moda dessa distribuição ocorre no intervalo ( 0 ; 10 ], com mais de 27 \% do total de participantes. 7.3.2.2 COMENTÁRIOS SOBRE AS RESPOSTAS À QUESTÃO DISCURSIVA 3 A questão abordou um tema importante para o Direito, visto que apresentou uma abordagem interdisciplinar, contemplando o Direito Ambiental, o Direito Civil e o Direito Penal.
O tema faz parte da formação básica do curso e pertence às suas Diretrizes Curriculares Nacionais.
Esperava -se do estudante um conhecimento básico para observar o caso concreto e explanar as consequências em diversos ramos do Direito.
Esperava -se que o estudante demonstrasse conhecimentos mínimos a respeito da responsabilização tanto da pessoa jurídica quanto de seus dirigentes e do Estado.
Eram aguardadas Respostas objetivas e claras dos estudantes.
O espaço disponível para a resposta dos três itens foi adequado e grande parte das respostas ocuparam em torno de 9 linhas.
Ao analisar as respostas, no item A, boa parte dos alunos respondeu de forma recorrente que a Companhia Vale deverá indenizar as famílias tanto pelo dano material quanto pelo dano moral, não exemplificando as naturezas dos danos e qual tipo de indenização ou pensão aplicar a cada caso.
No item B, observamos as respostas mais parecidas entre os alunos, citando de forma incisiva que os dirigentes da Companhia deveriam ser presos e reparar o dano ambiental, não trazendo as sanções à pessoa jurídica, a qual responderia isoladamente, independente do dolo praticado pelos dirigentes, raramente utilizando de jurisprudência.
O item C trouxe o maior índice de respostas incorretas, tendo a resposta recorrente de que o Estado fiscalizou a barragem de forma irresponsável e, por este motivo, a responsabilidade seria solidária.
Dessa forma, podemos concluir que o rendimento na Questão 3 foi baixo.
Com relação à linguagem utilizada pelos estudantes, as respostas estavam dentro do esperado para alunos que estão cumprindo o curso de Direito, utilizando de forma correta o vernáculo jurídico.
Com relação à letra A, observamos mais respostas preenchidas, porém muitos candidatos responderam o item somente indicando a indenização material e o dano moral, sem aprofundar e fundamentar o argumento.
Poucos conseguiram contemplar de forma completa os danos e as indenizações cabíveis às partes que sofreram com o rompimento da barragem.
Muitas respostas do item B não foram realizadas conforme o padrão de resposta.
Contemplando de forma sucinta que os dirigentes deveriam ser presos e obrigados a reparar o dano ambiental, não auferindo a abrangência completa que o item requeria, como abordar o Direito Penal, o Direito Ambiental e, ainda, especificar que tanto a pessoa jurídica quanto os dirigentes poderiam responder separadamente por suas condutas delituosas e sem a utilização de fundamentação e jurisprudência.
Muitos estudantes erraram o item C por completo, indicando, em sua resposta, que o Estado era responsável pela fiscalização da barragem e tão somente por isso teria responsabilidade ; por vezes respondiam que a natureza da responsabilidade era subsidiária, o que não está correto.
Ficou evidente que os estudantes tinham o mínimo de conhecimento das bases do dano moral, material e, ainda, sobre a aplicabilidade do Direito Ambiental, contudo, raras respostas utilizaram a jurisprudência como apoio em suas soluções do caso em comento.
A maior dificuldade foi de especificar e encontrar, de forma correta, a responsabilidade do Estado em caso de desastres naturais.
Foi identificado que a principal deficiência na formação dos estudantes estava relacionada à capacidade de elaborar respostas completas e fundamentadas, conforme espera -se de um operador de Direito, utilizando as leis e jurisprudências.
Ademais, houve a identificação correta do dano moral e material, sendo contemplado pela grande maioria, assim como a diferenciação entre pessoa física e jurídica para responder de forma distinta sobre seus atos delituosos. 7.3.2.3 ANÁLISE DA QUESTÃO DISCURSIVA 4 DO COMPONENTE ESPECÍFICO Na Tabela 7.19, constam as informações relativas à Questão 4 do conjunto de Questões do Componente de Conhecimento Específico.
Como se observa, a partir dos dados da Tabela 7.19, a Média geral do Brasil foi 24, sendo a maior Média registrada na região Sudeste ( 24,7 ), e a menor, na região Centro-Oeste ( 22,3 ).
Quanto à variabilidade das notas, o desvio-padrão geral do Brasil foi 19,9 ; sendo o maior desvio-padrão registrado na região Nordeste ( 20,3 ), e o menor, nas regiões Sul e Sudeste ( 19,7 ).
A nota Máxima ( 100,0 ) foi atingida por, pelo menos, um concluinte em todas as regiões.
A Mediana em todo o Brasil foi 25 ; o mesmo valor foi obtido nas regiões Sul e Sudeste.
Nas regiões Norte, Centro-Oeste e Nordeste a mediana foi de 22,5.
No Gráfico 7.12, está representada a distribuição de notas da Questão Discursiva 4, do Componente de Conhecimento Específico.
A distribuição mostrada no Gráfico 7.12, excetuando -se resultados nulos, tem moda no intervalo [ 0 ; 10 ] com mais de 17 \% do total de participantes.
Chama a atenção que mais de 25 \% dos participantes deixaram essa questão em branco. 7.3.2.4 COMENTÁRIOS SOBRE AS RESPOSTAS À QUESTÃO DISCURSIVA 4 A questão estava inserida na parte de Componente Específico e envolveu um estudo de caso relacionado a uma experiência real noticiada em diversos jornais.
A situação narrada abordou fatos atinentes às diversas áreas de atuação do Direito, com tema multidisciplinar envolvendo Direito Público e Privado.
Consoante ao Parecer CNE/CES n.º 146/2002, a questão tinha o intuito de avaliar diversas competências e habilidades que se objetivam com a formação do profissional da área, a citar a interpretação e a aplicação do Direito.
Esperava -se do estudante um conhecimento básico de conteúdos referentes às áreas de Direito Civil, Direito Constitucional, Direito do Trabalho e Direito Penal, inclusive de Direitos Humanos, prezando pela priorização de interdisciplinaridade indicada para as formações geral e técnico – jurídica, conforme Resolução CNE/CES n.º 5/2018.
Os fatos narrados pela questão puderam evidenciar violações aos Direitos Humanos e às leis nacionais, possibilitando avaliar a capacidade prático – profissional do estudante em aplicar os conhecimentos técnicos e apontar questões de responsabilização civil, penal e trabalhista pela violação aos Direitos Humanos.
Com relação ao questionamento sobre a responsabilidade trabalhista, era esperado que o estudante demonstrasse conhecimentos mínimos sobre proteção do trabalhador e da relação de trabalho, bem como sobre questões constitucionais atreladas à relação de trabalho em condições análogas à escravidão.
Com relação aos questionamentos sobre a responsabilidade penal e civil, era esperado que o estudante demonstrasse capacidade de associar as condutas apresentadas pelo caso e sistematicamente indicar as sanções penais conforme o tipo penal e os detalhes dos fatos, bem como indicar a responsabilidade de reparar os danos causados pela conduta ilícita apresentada.
De modo geral, as respostas podiam ser objetivas e capazes de demonstrar os conhecimentos necessários para que o resultado da avaliação constatasse a capacidade do estudante de observar e analisar o caso apresentado pela questão e, em seguida, problematizar os fatos, de modo a aplicar os conhecimentos técnicos – jurídicos com a devida contextualização aos elementos da questão.
A maioria dos estudantes respondeu conforme o esperado : identificar os fatos narrados e associá -los às nuances constitucional, civil, penal e trabalhista do Direito.
No entanto, uma parcela dos estudantes avaliados combinou as respostas dos três itens avaliados – análise de responsabilidade penal dos empregadores ; apresentação das consequências jurídicas na esfera trabalhista ; apresentação das consequências cíveis e constitucionais –, eventualmente criando um equívoco entre as responsabilidades.
Quanto ao item A, poucas respostas deixaram de apontar as condutas praticadas contra a vítima como condutas típicas da “ redução à condição análoga a de escravo ” – conforme o artigo 149 do Código Penal Brasileiro.
No entanto, grande parte dos estudantes deixou de indicar critérios agravantes para a imputação penal, como o fato de a vítima ter sofrido com as condutas típicas enquanto criança/adolescente.
Quanto ao item B, boa parte dos estudantes avaliados foi capaz de apontar a responsabilidade trabalhista dos empregadores em pagar as verbas devidas.
No entanto, na mesma proporção, uma parcela dos estudantes deixou de indicar explicitamente as verbas rescisórias devidas, bem como o necessário reconhecimento do vínculo empregatício.
Quanto ao item C, os estudantes mencionaram o direito da doméstica em reclamar por indenização por danos morais, mas deixaram de explicar e contextualizar a situação.
Poucos estudantes fundamentaram suas respostas acerca da responsabilidade civil, sem mencionar especificamente o Código Civil, nem a Constituição Federal, tampouco indicaram as características que apontavam a ilicitude dos atos e a violação dos direitos constitucionalmente protegidos.
Quanto ao item A, alguns estudantes tipificaram a conduta como “ cárcere privado ” e submeteram a resposta do item B à redução à condição análoga a de escravo.
Outros estudantes indicaram a conduta como sendo meramente “ trabalho forçado ”, ainda que este seja apenas um dos termos integrantes da tipificação penal da conduta que era esperada de se identificar pelo estudante.
Com relação ao item B, os estudantes deixaram de caracterizar as verbas trabalhistas como sendo de gênero rescisório, limitaram -se a especificar o pagamento de verbas trabalhistas de modo geral como consequência jurídica aos atos praticados pelos empregadores.
Em outras frequentes respostas, os estudantes justificavam as irresponsabilidades e as condutas ilícitas dos empregadores, sem especificamente apontar as consequências jurídicas, conforme pedia o comando do item.
Com relação ao item C, os estudantes raramente fizeram menção ao princípio constitucional da dignidade da pessoa humana, basilar para justificar a ilicitude das condutas praticadas pelos empregadores.
Geralmente, a responsabilidade apontada pelos estudantes se limitava aos danos morais.
Assim, poucos apresentaram respostas que abarcavam danos existenciais, bem como, raras vezes foi possível evidenciar uma contextualização com os elementos que configuram atos ilícitos.
Boa parte dos estudantes apontou “ dano moral ” como resposta, sem indicar como consequência jurídica a necessária reparação pelos danos causados e a possibilidade de a vítima reclamar pelos danos sofridos.
Os erros se relacionaram com a falta de fundamentação das respostas e de detalhes importantes que subsidiam as análises de responsabilidade penal e das consequências cível e trabalhista.
Com relação ao item A, parte dos estudantes deixou de complementar a análise de responsabilidade com o fato do crime ter sido praticado contra criança/adolescente.
Isso reflete a falta de profundidade em conhecimentos técnico-jurídicos em determinadas áreas do Direito.
Com relação aos itens B e C, os estudantes deixaram de detalhar aspectos cruciais quanto a classificações, princípios e tópicos introdutórios relacionados aos conhecimentos teóricos das áreas de Direito do Trabalho, Cível e Constitucional.
Ficou evidente, com a extensiva gama de conteúdos e detalhes abordados pela questão, que os estudantes possuem conhecimentos básicos, sem profundidade de especificidades em variados ramos do Direito.
Isto porque as respostas mais bem fundamentadas abarcavam um ou dois dos itens, deixando lacunas na argumentação dos demais aspectos solicitados pela questão.
A aplicação dos conhecimentos teóricos pelo profissional do Direito consiste no exemplificativo estudo de caso apresentado pela questão.
A prática jurídica que compõe obrigatoriamente o currículo dos cursos de Direito costuma ser o espaço de construção e consolidação das habilidades e competências práticas, justamente para a aplicação daquilo que é construído na jornada acadêmica dos estudantes.
A análise das respostas à questão aponta informações preocupantes sobre a capacidade prática dos profissionais egressos dos cursos de Direito. 7.3.2.5 ANÁLISE DA QUESTÃO DISCURSIVA 5 DO COMPONENTE ESPECÍFICO Na Tabela 7.20, constam as informações relativas à Questão 5 do conjunto do Componente de Conhecimento Específico.
Como revelam os dados da Tabela 7.20, o desempenho dos estudantes de todo o Brasil na Questão Discursiva 5 do Componente de Conhecimento Específico foi inferior ao de todas as demais questões discursivas.
A nota média dos estudantes de todo o Brasil foi 14,3.
A maior Média foi registrada na região Sul ( 14,7 ), enquanto a menor Média foi registrada na região Centro-Oeste ( 13,2 ).
Quanto à variabilidade das notas, o desvio-padrão dos estudantes do Brasil, como um todo, foram 15,5.
Enquanto o maior desvio foi encontrado na região Nordeste ( 16 ), o menor foi encontrado na região Sul ( 15,2 ).
A mediana para o Brasil e para todas as regiões foi 12,5, exceto a região Centro-Oeste e Norte que ficaram com mediana 10.
Para o conjunto de estudantes de Direito do Brasil, a nota máxima foi 100,0, obtida por, pelo menos, um estudante em cada região.
A nota mínima do Brasil foi zero, o mesmo valor encontrado em todas as regiões.
No Gráfico 7.13, é representada a distribuição das notas da Questão Discursiva 5 do Componente de Conhecimento Específico.
A distribuição mostrada no Gráfico 7.13, excetuando -se resultados nulos, tem moda no intervalo ( 0 ; 10 ], com 20 \% do total.
Nessa questão, pouco mais de 40 \% a deixaram sem resposta. 7.3.2.6 COMENTÁRIOS SOBRE AS RESPOSTAS À QUESTÃO DISCURSIVA 5 A questão aborda tema relacionado à atualidade, momento em que o mundo vivenciou uma situação de pandemia, sendo o sistema de saúde brasileiro bastante demandado e colocado à prova em um contexto de extrema necessidade, decorrente da propagação do vírus da Covid-19.
A questão caracteriza, ainda, situação que foi, de fato, evidenciada e vivenciada pelo sistema judiciário brasileiro, tendo recebido grande demanda de ações judiciais em que se questionou a atuação do Sistema Único de Saúde, o fornecimento de medicamentos, o tratamento e as vacinas. atuação do Sistema Único de Saúde, o fornecimento de medicamentos, o tratamento e as vacinas.
Importa salientar que a saúde, sendo um direito fundamental, é um tema que vem bem delineado pela Constituição Federal/88, que, ao instituir um federalismo solidário, impôs a competência comum da União, dos estados, do Distrito Federal e dos municípios para cuidar da saúde, nos termos do art. 23, II, da Carta Magna.
Trata -se de um tema elementar para os estudantes da área jurídica, sendo bem detalhado nos primeiros semestres do curso de Direito, formando o alicerce para o aprofundamento dos demais direitos que abrangem o tema.
Vale ressaltar que, nessa questão, esperava -se que o estudante dominasse as informações básicas e substanciais para o conhecimento sobre o direito fundamental à saúde, a distribuição de competências para o alcance das respectivas responsabilidades, além do mecanismo judicial a ser utilizado por aqueles que tivessem o seu direito lesado ou negado.
Apesar de tratar -se de questionamento com demanda para respostas discursivas, ou seja, prova subjetiva, as respostas precisavam ser claras e diretas, na medida em que os questionamentos foram de conhecimentos básicos, não se tratando de posicionamentos divergentes ou mesmo que exigissem formação de juízo de valor.
Foi fornecido um espaço suficiente para resposta, delimitado em 15 linhas, porém boa parte dos estudantes não \textbf{ultrapassou} cinco linhas para resposta, além de ter respondido a apenas um dos itens questionados.
Observou -se, ainda, que muitos candidatos respondiam aos itens em forma de redação corrida, sem topicalização das respostas.
Observou -se que os estudantes, de modo geral, não conseguiram reproduzir um conhecimento mais detalhado sobre o tema, em especial sobre a repartição das competências no que se refere à saúde.
Em relação ao posicionamento do STF, uma pequena parte explanou essa resposta, mas sem demonstrar grande segurança e riqueza de detalhes/argumentos.
Já em relação às medidas judiciais, foi possível identificar mais compreensão do assunto, apesar das respostas serem em regra desprovidas de fundamentação.
Importa salientar que a maioria dos estudantes não demonstrou domínio ou intimidade com a linguagem técnico – jurídica, de modo que a linguagem comum foi utilizada pela maioria.
A questão era composta por duas perguntas.
Com relação ao item A da questão, muitos estudantes confundiram o sistema de repartição de competências, demonstrando total desconhecimento acerca da competência concorrente entre União, estados, Distrito Federal e municípios.
Em relação ao entendimento do STF, de que todos os entes federativos possuem competência para assegurar a efetividade e plenitude da saúde pública, poucos estudantes citaram e fundamentaram essa posição.
Boa parte não abordou o entendimento do STF, limitando -se a falar sobre competência comum entre os entes, sem quaisquer esclarecimentos.
Em relação ao item B, muitos candidatos responderam de forma bem direta que a medida judicial cabível seria o mandado de segurança, sem maiores esclarecimentos ou fundamentação mais detalhada.
Uma parte dos estudantes também identificou a ação ordinária como caminho, com pedido de tutela de urgência.
Dentre esses, poucos chegaram a esclarecer a presença da probabilidade do direito e do perigo de dano ou risco ao resultado útil do processo. 7.3.2.7 CONSIDERAÇÕES FINAIS SOBRE AS QUESTÕES DISCURSIVAS DO COMPONENTE DE CONHECIMENTO ESPECÍFICO Ficou evidente que os estudantes precisam treinar tanto a habilidade de escrita, tão exigida de um profissional da área jurídica, quanto a aplicação do raciocínio jurídico adquirido ao longo do curso.
Observa -se que os estudantes possuem acesso ao conhecimento, mas, no momento de aplica -lo, a insegurança proveniente da pouca prática termina prevalecendo e ocasionando um resultado insuficiente em fundamentos jurídicos, em linguagem técnica e em soluções jurídicas decorrentes de um conhecimento mais abrangente adquirido ao longo do curso.
Identificou -se que os estudantes, apesar de terem tido contato com o tema objeto da questão, possuem dificuldade em aplicar o conhecimento jurídico transmitido pelas instituições.
Percebe -se uma dificuldade na escrita técnico – jurídica e, também, na aplicação do raciocínio jurídico necessário ao deslinde do caso apresentado.
As dificuldades identificadas talvez estejam relacionadas à pouca prática desses candidatos no que se refere à habilidade de responder questões subjetivas, que exigem textos dissertativos.
Percebe -se que boa parte dos estudantes não possuem qualidade de vocabulário, nem intimidade com a escrita gramaticalmente correta, o que ocasiona a produção de textos pobres em argumentação, com linguagem informal e, muitas vezes, repetitiva, bem como textos extremamente curtos e desprovidos de argumentos básicos.

% matched lemmas: brasileira, observação, ultrapassar
\end{textsample}
